\section{背景与问题定义}
\label{sec:background}

\subsection{MiniOneRec tokenizer pipeline}

设每个 item $i$ 对应一个语义 embedding $\mathbf{x}_i \in \R^d$。MiniOneRec 首先在这些 item embedding 上训练一个 \rqvae{}，然后将每个 item 编码成固定长度的 semantic ID，记为 $\sid(i) = [z_i^{(1)}, z_i^{(2)}, z_i^{(3)}]$ \citep{kong2025minionerec,lee2022rqvae}。在 \texttt{sid-train} 阶段，模型学习对语义空间进行残差量化；在 \texttt{sid-generate} 阶段，系统再对 colliding items 启用 Sinkhorn-based reassignment 进行多轮重编码，从而显著降低最终导出的 \texttt{index.json} collision。

这一点很重要。tokenizer 的质量不能只看 \texttt{sid-train} 阶段的 raw collision，因为真实 pipeline 最终使用的是 \texttt{sid-generate} 之后的 index。因此，公平的 tokenizer 比较必须落在最终导出的 SID 上，而不是只看训练器中间指标。

\subsection{为什么当前更像 local ambiguity 问题}

我们对当前 semantic baseline 的分析持续指向同一个失败模式：许多错误发生在模型已经进入正确语义区域之后。更具体地说，第一层甚至前两层 prefix 往往是正确的，但最后一层 leaf token 仍然出错。于是，当前剩余难点更像是 \emph{语义合理前缀内部的局部歧义}。

为此，我们把两级前缀 $(z_i^{(1)}, z_i^{(2)})$ 相同的 item 集合称作 same-$l_2$ bucket。如果某个 prefix 下仍然挂着很多 leaf，那么即使上层 routing 已经正确，生成器在最后一步仍然要面对困难的局部判别。基于这一观察，本文除了报告 collision 之外，也把 same-$l_2$ 的拥挤程度和给定 prefix 条件下第三级 code 的条件熵当作 tokenizer 健康度指标。

\subsection{公平比较要求对齐上游训练制度}

在比较 graph-aware tokenizer 之前，我们首先遇到了一个 provenance 问题。用当前仓库默认配置直接 fresh rerun \texttt{sid-train -> sid-generate}，会得到非常高的 collision，这与仓库中已经存在的 Industrial 正式 index 明显不符。通过逐步对比本地代码与上游 MiniOneRec 官方仓库，并重新用上游训练制度跑一遍 Industrial tokenizer，我们确认问题主要来自超参数而不是实现损坏。

\begin{table}[t]
  \centering
  \caption{Industrial 上 tokenizer provenance 对齐结果。只有使用上游风格 \rqvae{} 训练制度，才能复现当前仓库正式 tokenizer 的低 collision 质量。}
  \label{tab:provenance}
  \begin{tabular}{lcccc}
    \toprule
    设定 & Epochs & Batch size & 最佳 train collision & 最终 index collision \\
    \midrule
    本地默认 fresh rerun & 10 & 256 & 0.9989 & 0.9940 \\
    上游风格 fresh rerun & 10000 & 20480 & 0.0993 & 0.00434 \\
    仓库正式 Industrial index & -- & -- & -- & 0.00434 \\
    \bottomrule
  \end{tabular}
\end{table}

\Cref{tab:provenance} 是本文后续实验的前提。它说明，任何 tokenizer-side 方法比较都必须建立在 upstream-aligned 的训练制度之上，否则连 baseline 都是不公平的。因此，本文所有训练时的 MGR-SID 比较都采用与上游 MiniOneRec 一致的 \rqvae{} 训练设置。

\subsection{问题定义}

基于这些事实，本文关注的问题不再是泛泛的“是否要把 collaboration 融进 tokenizer”，而是：
\begin{quote}
在不破坏 semantic hierarchy 的前提下，graph-structured collaborative information 应该怎样以层级感知的方式进入不同 SID level，从而降低 local ambiguity？
\end{quote}

当前 MGR-SID v1 对这一问题的回答是保守而清晰的：保留语义 \rqvae{} backbone，不去替换 semantic tokenizer，而是把图结构作为 level-specific structural supervision 注入到 tokenizer 学习过程中。
